% =====================================================================
\chapter[Structural Convergence]{Structural Convergence: Over-Constrained Solution Sets and Engineering Attractors}\label{ch:convergence}
% =====================================================================

This chapter addresses Problem \ref{prob:low}. High-level structure does converge, but by a
different mechanism from the low level, and claims of different strength must be stated
separately; the strength of convergence depends on domain binding, and its carrier is
relations rather than components. The chapter closes with a falsifiable verification design.

\section{Three Convergence Claims of Different Strength}\label{sec:three-claims}

\begin{hypothesis}[The weak claim: recurrence of local patterns]\label{hyp:weak}
Software in different domains repeatedly adopts the same set of local structural solutions:
layering and staging, separation of mechanism from policy, separation of control from data,
events and state machines, separation of declarative description from execution engine,
separation of front end from back end, and resource ownership and lifecycle management.
What matters about these structures is not that they share a name, but that they share a
\emph{structural motivation}.
\end{hypothesis}

The weak claim has substantial engineering experience behind it, and it is also the weakest:
it is almost impossible to falsify, and therefore has limited explanatory power.

\begin{hypothesis}[The medium claim: structural similarity after abstraction]\label{hyp:medium}
For large software systems that come from different domains, are developed by different
teams and carry different implementation details, once differences of implementation
language, component naming and domain semantics have been eliminated, their high-level
architectures exhibit a structural similarity \emph{significantly above a random baseline}
in some definable abstract space.
\end{hypothesis}

The medium claim is the core claim, and the operationalisable one. It can be operationalised
(\S\ref{sec:baseline}), and whether it holds depends on no metaphysical assumption.

\begin{hypothesis}[The strong claim: single-point convergence]\label{hyp:strong}
Almost all large software that is normally biased eventually converges on one and the same
high-level architecture.
\end{hypothesis}

The strong claim is unsupported by the evidence and is probably too strong. Its principal
defect is that it conflates ``there exists an attractor'' with ``there is only one
attractor''.

This paper takes Hypothesis \ref{hyp:medium} (the medium claim) as its core claim, treats
Hypothesis \ref{hyp:weak} (the weak claim) as its degenerate form, and rejects
Hypothesis \ref{hyp:strong} (the strong claim) as an overstatement.

\section{Constraint Systems and Solution-Set Geometry}\label{sec:overconstrain}

\begin{definition}[Constraint system and solution set]\label{def:constraint}
Let the design space be a $d$-dimensional manifold $\mathcal{D}$. A design problem is
defined by $c$ independent constraints
$g_1, \dots, g_c : \mathcal{D} \to \mathbb{R}$, and its \term{solution set} is
$\mathcal{S} = \{x \in \mathcal{D} : g_j(x) = 0,\; j = 1..c\}$.
\end{definition}

\begin{proposition}[Over- and under-constraint]\label{prop:overconstrain}
Under the notation of Definition \ref{def:constraint}, if the constraints are mutually
independent and transversal, then
\begin{enumerate}[label=(\alph*)]
\item when $c \ge d$, $\mathcal{S}$ is generically a finite set (or empty);
\item when $c < d$, $\mathcal{S}$ is generically a submanifold of dimension
      $\ge d - c$;
\item when $c \ll d$, the dimension of $\mathcal{S}$ tends to $d$, that is, the solution
      set is almost unconstrained.
\end{enumerate}
\end{proposition}

\begin{proof}
Under transversality, each independent constraint lowers the codimension by $1$. By the
transversality theorem (Thom), codimension is additive:
$\mathrm{codim}\,\mathcal{S} = c$, hence $\dim \mathcal{S} = d - c$ (when $d-c \ge 0$);
when $c \ge d$ the dimension is $0$, that is, a discrete point set; when $c < d$ the
dimension is positive, that is, a submanifold.
\end{proof}

Proposition \ref{prop:overconstrain} gives a clean explanation of ``unique below, manifold
above'', but the wording must be handled with care. It does \emph{not} say that the low
level really is formalised as a constraint system; it says that the effective dimension of
the low-level design space is small and its effective constraints numerous, so that the
solution set is degenerate in the geometric sense.

\begin{example}[The low level]\label{ex:abi}
Consider defining a function calling convention for a fixed instruction-set architecture.
The design freedoms include where arguments are passed (registers or stack), who is
responsible for saving registers, stack frame layout, where the return address is stored,
how variadic arguments are passed, alignment requirements, and compatibility with existing
binaries, debuggers, exception unwind tables, the system-call interface and the dynamic
linker. The great majority of these constraints are not negotiable preferences but are
fixed by \emph{objects that already exist}: existing binaries must keep working, the
debugger must be able to walk the stack, exceptions must unwind, and the kernel must be
able to receive system calls. Hence $c \ge d$, and the solution set degenerates to very few
feasible points, in practice often to a single point.
\end{example}

\begin{example}[The high level]\label{ex:arch}
Consider choosing an architecture for a task-processing system. The design freedoms
include: function calls within a monolith, multi-process message passing, microservice RPC,
event-driven asynchronous processing, the actor model, distributed workflow engines, and
their continuous parameters along the following dimensions: deployment granularity, state
ownership, consistency model, synchronous or asynchronous, failure boundaries, and data
partitioning. Here the problem-level constraints are markedly fewer than the degrees of
freedom, and $\mathcal{S}$ is a high-dimensional manifold.
\end{example}

Why, then, does the high level still converge? Because there is a \emph{second set of
constraints}, which come not from the problem but from the human organisations that build
and maintain the system:

\begin{principle}[Compression of the solution set by organisational constraints]\label{pr:org}
The under-constrained high-dimensional solution manifold at the problem level is compressed
by organisational constraints into a few sharp peaks. These constraints include: the cost of
communication between teams, the ceiling on cognitive load, the limits of deployment and
operations capacity, the market for technical talent, the sunk cost of an existing
technology stack, and regulatory and compliance requirements.
\end{principle}

One corollary follows: \emph{high-level convergence is not towards engineering optimality
but towards structures navigable by a bounded organisation}. The strength law below makes
this precise.

\section{Mechanisms of Convergence}\label{sec:mechanisms}

\begin{enumerate}[label=(M\arabic*)]
\item \term{Complexity constraint}. Let a system have $N$ components. If dependencies are
      left unorganised, the number of potential directed dependencies reaches $N(N-1)$, and
      the undirected count $N(N-1)/2$. Maintainability therefore requires structure to be
      imposed on dependencies: modularisation, layering, explicit interfaces, constraints on
      dependency direction, subsystem boundaries. These techniques are not identical, but
      all of them solve one and the same problem: \emph{limiting the propagation of
      complexity}. The classical argument for modularisation and information hiding is
      \cite{parnas1972}.
\item \term{Information-flow constraint}. A large class of software can be abstracted as
      ``input $\to$ processing $\to$ state change $\to$ output'' (with feedback, concurrent
      events, external side effects and continuously running state machines superimposed).
      Where different duties have different rates of change, error modes and resource
      requirements, separating them yields a concrete benefit: parsing from execution,
      business rules from data access, pure computation from side effects, and the control
      plane from the data plane.
\item \term{State-management constraint}. A system must answer: who owns a given state? who
      may modify it? when is it created and destroyed? how do several components observe its
      changes? how is concurrent access coordinated? how is consistency restored on failure?
      These questions are highly general, and therefore repeatedly produce similar
      structures: transaction and lock management, process and handle management, the state
      owner and event loop of a GUI, replicated state machines, entity lifecycle and system
      update.
\item \term{Cognitive constraint}. Large software is constrained not only by machines but by
      the limits of developer comprehension. If implicit relations hold between all
      components, then even a system that runs in the computational sense leaves its
      developers unable to predict the effect of a change. Humans therefore tend to build
      hierarchies, namespaces, modules, explicit interfaces, composable abstractions and
      locally comprehensible state. \emph{Software architecture organises not only
      computation but comprehensibility under bounded cognition.}
\item \term{Evolutionary constraint}. Large software is not designed in one pass; it passes
      through initial implementation, feature growth, performance optimisation, fault repair,
      team handover, module refactoring and compatibility maintenance. Some decisions reveal
      their cost only as evolution proceeds, which produces a selection effect. Evolution in
      practice does not guarantee optimality: legacy baggage, compatibility requirements,
      organisational structure and sunk cost can all keep an undesirable architecture alive
      for a long time.
\end{enumerate}

Mechanisms (M1)--(M3) come from the objective structure of the problem; mechanisms (M4)--(M5)
come from humans and history. This distinction corresponds to the strength law of
\S\ref{sec:strength}: the first three mechanisms produce cross-domain similarity, the last
two produce within-domain similarity and path dependence.

The classical statements of mechanisms (M4) and (M5) are decomposition into levels under
cognitive load \cite{simon1962} and the determination of system structure by communication
structure \cite{conway1968}. The oscillation between a ``no final state'' view of
architecture and multiple stable forms is treated in \cite{nygard2018}; the practice by
which organisational shape drives architecture back through team topologies appears in
bounded contexts and domain-driven design \cite{evans2003}.

\section{Convergence Strength and Domain Binding}\label{sec:strength}

\begin{definition}[Domain binding]\label{def:dom}
Let the semantics of level $\Lambda$ depend on a set of domain axioms
$\mathrm{Ax}(\Lambda)$ (for example, ``an account balance may be negative'' or ``a frame
must be presented within 16.6\,ms''). Define the \term{domain binding} of $\Lambda$ as
\begin{equation}\label{eq:dom}
\mathrm{Dom}(\Lambda) \;=\; \frac{\bigl|\mathrm{Ax}(\Lambda)\bigr|}
{\bigl|\mathrm{Ax}(\Lambda)\bigr| + \bigl|\mathrm{Ax}_{\text{general}}(\Lambda)\bigr|} ,
\end{equation}
that is, the fraction of the level's axioms that are domain-specific. $\mathrm{Dom} = 0$
means that the level is determined entirely by general axioms (such as memory locality or
the cost of concurrency coordination); $\mathrm{Dom} = 1$ means that the level is entirely
domain convention.
\end{definition}

\begin{conjecture}[The strength law of convergence]\label{conj:strength}
Let $\mathrm{Conv}(\Lambda)$ be the structural similarity of level $\Lambda$ under
cross-system comparison (\S\ref{sec:baseline} gives the operationalisation). Then
$\mathrm{Conv}$ is monotonically decreasing in $\mathrm{Dom}$:
\begin{equation}\label{eq:convlaw}
\Lambda_1 \preceq \Lambda_2 \;\Longrightarrow\;
\mathrm{Conv}(\Lambda_1) \;\ge\; \mathrm{Conv}(\Lambda_2)
\qquad \text{whenever} \qquad
\mathrm{Dom}(\Lambda_1) \le \mathrm{Dom}(\Lambda_2).
\end{equation}
\end{conjecture}

Conjecture \ref{conj:strength} has observable qualitative support: the closer to the
physical and network layers, the stronger the convergence (the eight fallacies of
distributed computing are hard constraints with no room for negotiation); the closer to the
business domain, the weaker the convergence (an insurance core system and a recommender
system are almost never isomorphic in their domain models); the closer to organisational
structure, the more it is governed by communication structure; the closer to the maturity
phase of a technology, the more its form converges; and the closer to core and high-risk
modules, the more conservatively it converges.

This paper does not call Equation \eqref{eq:convlaw} a law, because it depends on the
definability of $\mathrm{Dom}$ and the measurability of $\mathrm{Conv}$, both of which must
be established experimentally. Its value is to convert a vague directional judgement into a
falsifiable monotonicity hypothesis.

\section{Relations Rather Than Components}\label{sec:relations-converge}

\begin{principle}[Relations first]\label{pr:relations}
Cross-system structural commonality resides mainly in \emph{relations} rather than in
\emph{components}. Different systems need not have the same modules, but they may have
similar relational graphs:
\begin{itemize}
\item which parts decide and which parts execute;
\item which parts own state and which parts may only read it;
\item which parts define stable contracts and which parts are allowed to change;
\item which parts coordinate globally and which parts work independently;
\item which dependencies may cross a boundary and which must be forbidden.
\end{itemize}
\end{principle}

In the language of Definitions \ref{def:relgraph} and \ref{def:skeleton}: what must be
sought is not a universal architectural template but a set of structural relations that
recur across different systems, together with the \emph{conditions under which they hold}.
This is the same line of thought as the six functions of Chapter \ref{ch:infrastructure},
viewed from two sides: the six functions supply coordinates for \emph{management duties},
the relational graph supplies coordinates \emph{between duties}.

\begin{example}[Abstract skeletons of a database and a browser]\label{ex:db-browser}
A database engine comprises an SQL parser, a query optimiser, an execution engine, a buffer
pool, a transaction manager and a storage engine; a browser comprises HTML/CSS/JavaScript
parsers, a style and layout engine, a script execution engine, a rendering engine, a network
subsystem and a cache persistence mechanism. Their module names, internal algorithms and
concrete duties all differ. Yet under a suitable skeleton map $\phi$, both present as: an
\emph{input and interface layer} (receiving requests, parsing input, establishing an
execution context); an \emph{interpretation and planning layer} (converting external intent
into an internally executable representation); an \emph{execution and coordination layer}
(scheduling work, coordinating subsystems, managing the execution process); and a
\emph{state and resource layer} (managing data, caches, lifecycle and persistence).
\end{example}

The value of Example \ref{ex:db-browser} is limited: this skeleton is already broad enough
to describe almost any system. An additional requirement on the abstraction level is
therefore needed:

\begin{principle}[Two-way constraint on the abstraction level]\label{pr:abstraction-level}
An abstraction level worth studying must satisfy two directions at once: \emph{upwards}, it
must be able to identify structure common across systems; \emph{downwards}, it must retain
enough discriminating power that different systems remain distinguishable at that level.
The more abstract the level, the more implementation differences are eliminated, but the
more easily the description loses discriminating power; there is a proper middle level.
\end{principle}

The coordinates of the six functions in Chapter \ref{ch:infrastructure} are designed
precisely to satisfy this two-way requirement: they are narrower than
``input--processing--output'' (and therefore discriminating) and broader than any concrete
module list (and therefore comparable across domains).

\section{Verification Design}\label{sec:baseline}

For the medium claim (Hypothesis \ref{hyp:medium}) to be a scientific claim, it must be able
to answer ``compared with what''. A complete test design has five steps.

\paragraph{Step 1: Build the architecture graph.}
For each system, build a graph $\mathcal{G} = (V, E)$ whose nodes are components or duties
and whose edges are calls, dependencies, control flow or data flow. By
Definition \ref{def:relgraph}, \emph{relations of different types must be recorded
separately}, so that entirely different relations are not conflated.

\paragraph{Step 2: Define the abstraction map.}
Map concrete components to more general categories, for example: input and parsing, planning
and decision, execution and coordination, state and persistence, communication and
boundary, resource and lifecycle management. These categories must not be defined so broadly
in advance that any system will appear similar.

\paragraph{Step 3: Choose discriminating structural metrics.}
Candidate metrics include: the degree of modularity of the dependency graph; the depth of
the hierarchy; the proportion of cross-module dependencies; the concentration of state
ownership; the degree of coupling between control flow and data flow; the network centrality
of core components; and the direction of dependencies between duties.

\paragraph{Step 4: Construct a sound control group.}
It is not enough to compare a few well-known systems and find that all of them are modular.
The observed similarity must be compared against an appropriate random baseline, controlling
for system size, domain, platform, architectural style and historical technology stack.

\paragraph{Step 5: Distinguish genuine convergence from common origin.}
Two systems may adopt the same architectural pattern because they independently solved
similar problems; because they used the same framework; because they were shaped by the same
standard; because they drew on the same design literature; or because they inherited the
same historical architecture. \emph{Only after these factors are controlled is there better
reason to speak of independent structural convergence.}

Chapter \ref{ch:program} writes this design out as Hypotheses $H_1$ and $H_2$.

\section{Counterexamples and Limits}\label{sec:limits}

The cases that invalidate or degrade the convergence claim must be listed explicitly:

\begin{enumerate}[label=(\alph*)]
\item \term{Path dependence}. Once a structure is locked in by an ecosystem, a toolchain or
      a standard, it may persist for a long time even if suboptimal. The strength of the
      lock-in is proportional to the cost of migration and independent of the quality of the
      structure locked in.
\item \term{Phase transition in size}. At the two sizes of $10^{2}$ and $10^{9}$ entities,
      the optimal structure is qualitatively different: the consistency model, the sharding
      strategy and the failure model all change. There is therefore no such thing as ``the
      architecture of a large system'', only ``an architecture within some range of sizes''.
      This directly limits any statement that ``all large systems converge''.
\item \term{Organisational politics}. The determination of system structure by communication
      structure (Conway's law) means that even when the technical problem is identical,
      different organisations will produce different system boundaries. This is both part of
      the reason for convergence (when organisational shapes converge) and part of the reason
      for divergence (when organisational shapes differ).
\item \term{Strongly coupled systems}. When the internal coupling of a system cannot be
      localised (for example, a tightly coupled real-time control system or a monolithic
      kernel), module boundaries may be forced rather than chosen, and comparing
      ``architecture'' itself becomes less meaningful.
\item \term{Arbitrariness in the choice of abstraction level}. The choice of metrics in
      Steps 3 and 4 itself contains researcher degrees of freedom. Different duty
      classification systems yield different similarity conclusions. This is the most fragile
      methodological link in this line of work and must be mitigated by pre-registering the
      metrics and by robustness tests across several classification systems.
\end{enumerate}

The falsifiable part of this chapter is Hypothesis \ref{hyp:medium} and
Conjecture \ref{conj:strength}: if, after controlling for size, domain, platform and
framework, cross-system structural similarity is not significantly above a random baseline,
or if similarity is not monotone in $\mathrm{Dom}$, then the central claim of this chapter is
weakened (Hypotheses $H_1$ and $H_2$, Chapter \ref{ch:program}). The weak claim
\ref{hyp:weak} and Principle \ref{pr:relations} are unaffected by this, because they are
interpretive frameworks rather than empirical assertions.
